\documentclass[10pt]{article}
\usepackage[margin=0.78in]{geometry}
\usepackage{amsmath,amssymb,booktabs,array}
\usepackage[T1]{fontenc}
\usepackage{lmodern}
\usepackage[colorlinks=true,linkcolor=blue,urlcolor=blue]{hyperref}

\title{TPC-389: Long-Horizon Stress Test for a Frozen Count Slope}
\author{Liang Wang\\
School of Mathematics and Statistics, Huazhong University\\
of Science and Technology (HUST), Wuhan, China}
\date{September 5, 2026}

\begin{document}
\maketitle

\begin{abstract}
TPC-388 reported a finite cross-family transfer of a logarithmic count slope.
TPC-389 asks whether that interface survives one longer count interval.  We
freeze the TPC-388 slopes and, on a third coordinate-disjoint family, use
three origins at counts $768$ and $1024$ for calibration and two later origins
at count $1280$ for holdout.  The panel has 256 rows and 32 cells.  Anchored
parent, same-family local, and recursive parent forecasts all pass a
predeclared 3\% finite ratio cap.  The largest errors are
$0.0176155841$, $0.0119975160$, and $0.0299499406$, respectively.  This is a
finite stress certificate: the inherited spectral diagnostic still fails on
64 of 256 rows, and no arithmetic or asymptotic conclusion follows.
\end{abstract}

\section{Question and claim boundary}

TPC-388 showed that 32 slopes learned on an earlier finite count ladder
survived transfer to a fresh origin family through $N=1024$.  The next minimal
question is horizon stress: does the same frozen interface remain predictive
from $N=1024$ to $N=1280$?  We choose a third, coordinate-disjoint affine
family
\[
 a_j=2800001+401j,\qquad 0\leq j<41,
\]
and freeze indices $0,10,20,30,40$.  The first three origins are calibration
origins and the last two are holdout origins.  All roles and the parent
interface are fixed before current responses are read.

The claim firewall is
\[
 \texttt{ARITHMETIC\_ADVANCE=NO},\qquad
 \texttt{FIXED\_POWER\_CREDIT=0},\qquad
 \texttt{FULL\_GATE\_B=OPEN}.
\]
The Session's official evaluator files are absent in this checkout; the local
Bridge-B is fail-closed repository evidence only.

\section{Finite proxy}

For $p\in(Q,2Q]$ and $H=66$, define
\begin{align*}
 K_p(u,v)={}&p(p/Q)^2\frac{H^2}{H^2+(u-v)^2}
 \left({\bf 1}_{p\mid u-v}-\frac{1}{p-1}\right)\\
 &\quad\cdot{\bf 1}_{u\ne v}{\bf 1}_{p\nmid u}{\bf 1}_{p\nmid v}.
\end{align*}
The row geometry is the finite square energy
$G(u)=\sum_{p}\sum_{v\in I}K_p(u,v)^2$.  For each of four sign laws
$\ell$, the matrix is $M_\ell(u,v)=\sum_p s_\ell(p)K_p(u,v)$.  We report
local-diagonal normalization
$M_\ell(u,v)/\sqrt{G(u)G(v)}$ and a pooled scalar normalization based only on
calibration-origin geometry.  The fixed band keeps block distance at most
three; the full-relative band keeps all block pairs.  We use
$Q\in\{2048,8192\}$ and block length 128.

\section{Forecast interface}

For a cell, let $S_N$ be the mean band spectral diagnostic over the relevant
origins.  The parent exponent $\alpha_{\rm P}$ is read from the hashed
TPC-388 certificate and is not refit.  The local exponent is fitted only from
the two current calibration means:
\[
 \alpha_{\rm L}=
 \frac{\log(S_{1024}/S_{768})}{\log(1024/768)}.
\]
The anchored forecasts and the recursive forecast are
\[
 \widehat S_{1280}^{\rm P}
   =S_{1024}(1280/1024)^{\alpha_{\rm P}},\qquad
 \widehat S_{1280}^{\rm L}
   =S_{1024}(1280/1024)^{\alpha_{\rm L}},
\]
\[
 \widehat S_{1280}^{\rm rec}
   =S_{768}(1280/768)^{\alpha_{\rm P}}.
\]
The holdout ratio is $S_{1280}/\widehat S_{1280}-1$; the finite pass cap is
$|\text{ratio}|\leq0.03$.  No ratio is interpreted outside this declared
finite panel.

\section{Certification and results}

The canonical certificate contains 256 rows and 32 cells.  The producer
accumulates prime shells in ascending order; an independent implementation
rebuilds them in descending order and recomputes row values and all forecast
cells.  A rational 13-point $Q=8$ anchor at
$[2800001,2800014)$ verifies positive geometry and symmetry for all four laws.
A 25-mutation stress test rejects altered roles, hashes, summaries, and
firewall values.

\begin{table}[t]
\centering
\caption{Finite forecast and stability census.}
\small
\begin{tabular}{@{}l r@{}}
\toprule
quantity & result\\
\midrule
parent anchored pass & 32/32\\
local-control pass & 32/32\\
recursive parent pass & 32/32\\
maximum parent error & 0.0176155841\\
maximum local error & 0.0119975160\\
maximum recursive error & 0.0299499406\\
stable cells ($N=768,1024,1280$) & 24/32,\ 27/32,\ 24/32\\
spectral failures / Schur failures & 64/256,\ 0/256\\
\bottomrule
\end{tabular}
\end{table}

The strongest anchored deviation is the fixed-three-block, pooled, all-plus
cell at $Q=8192$.  The strongest recursive deviation is fixed-three-block,
pooled, half-split at $Q=2048$.  The latter is close to, but still inside,
the finite cap.  This proximity is a stress observation, not a margin for an
asymptotic theorem.  The 64 spectral failures occur while all Schur diagnostics
remain below their cap, so the finite forecast success does not repair the
operator-norm obstruction.

\section{Conclusion and next clue}

TPC-389 supplies one independent finite progress point: a frozen slope remains
predictive over a longer count horizon on a third coordinate family, under
anchored, local, and recursive tests.  The reusable structure is
\[
 \text{hashed parent interface}\longrightarrow
 \text{fresh calibration/holdout ladder}\longrightarrow
 \text{three forecast audits}.
\]
The strongest obstruction is the growing spectral-cap failure census, together
with the recursive error being near the declared boundary.  The open theorem
is a source-valid, origin/count-uniform growing operator and slope control.
The next clue is
\[
 \texttt{ROUND2\_CLUE=TEST\_C1\_RECURSIVE\_SLOPE\_COMPOSITION}.
\]
Arithmetic reassembly, Route-A/Route-B closure, and a twin-prime theorem remain
open.

\section*{Reproduction}

All source, certificate, proof, experiment, and Bridge-B files are stored under
the project directory:

\url{papers/tpc-389-c1-long-horizon-slope-stress/}

Run the ordinary and optimized producer, independent checker, mutation stress,
and local Bridge-B commands listed in the project README.  The release uses
\texttt{paper/main.pdf} and \texttt{paper/paper.pdf} as byte-identical copies.

\end{document}
